%% ============================================================
%% RELATED WORK -- v2 merge, drop-in replacement for the
%% <<KEEP:rw-milp-topology>> <<KEEP:rw-thermohydraulic>> <<KEEP:rw-positioning>>
%% markers under \subsection{Related work} in the Introduction.
%%
%% Sources: v1 lines 317-500, remapped to the v2 taxonomy (v1 L1 = \CP,
%% v1 L3 = \Lone -- see LEVEL_CROSSWALK_v1_to_v2.md) and re-argued toward the
%% TWO gaps this revision claims (loss-vs-topology confound; decisions vs
%% objective values) rather than v1's single "topology vs physics" gap.
%%
%% Editor compliance: every lumped \cite{a,b,c} from v1 reduced to the single
%% most relevant reference. v1 lumps removed: {Fang2016,Guelpa2019review},
%% {Connolly2014,Brown2018,Lund2017}, {Wirtz2020,Mertz2016,Buffa2019},
%% {Vesterlund2017,Schweiger2017,vanderHeijde2017,Blommaert2020},
%% {Wirtz2020,Buffa2019}, {Schweiger2017_p2h,Averfalk2017,David2017},
%% {Bordin2016,Mertz2016}, {vanderHeijde2017,Giraud2015}.
%%
%% ONE OPEN ITEM flagged inline: %% <<CITE-NEEDED:vss>>
%% ============================================================

Two strands of the district-heating literature bear on this paper: how finely a
dispatch model resolves the network, and how finely it resolves the network's physics.
Neither strand has been evaluated by the quantity that matters for a scheduling model:
the cost of the decisions it produces.

Mixed-integer linear programming dominates operational district-heating
optimisation~\cite{Fang2016}, and the general-purpose frameworks in which such studies are
built (oemof~\cite{Hilpert2018_oemof}, PyPSA~\cite{Brown2018_pypsa},
FINE~\cite{Welder2018_FINE}, Calliope~\cite{Pfenninger2018_calliope},
urbs~\cite{Dorfner2016_urbs}) are typically applied with aggregated or component-based
thermal networks: losses time-invariant or absent, no hydraulic coupling, topology fixed at
one level of detail, at what this paper calls the copperplate level $\CP$; detailed
thermo-hydraulic coupling requires additional formulations. The cost of that simplification
is not quantified.

Between the copperplate and the full graph lies a continuous spectrum of spatial
aggregation. \citet{Larsen2004} compare aggregation methods on a 23-substation network
and show that pipe count can be reduced by roughly an order of magnitude before
production and return-temperature errors become significant; \citet{Falay2020} carry the
approach to dynamic simulation of several thousand substations. Both establish
\emph{which} aggregations preserve simulation fidelity, and neither quantifies the
resulting bias in dispatch-optimisation cost. The same emphasis persists in the most
recent work: \citet{Vrain2025} formalise an aggregation method with an explicit measure of
the error it induces in large-scale multi-energy \emph{simulation}, and \citet{Cassetti2025}
show how spatial resolution shapes the outcome of regional heat-network \emph{planning}.
Both confirm that aggregation error is resolution-dependent, yet both judge simulated or
planned quantities rather than the bias a coarse network imparts to the \emph{operational
dispatch decision}, and neither isolates thermal loss from spatial topology. Zone-based models aggregate demand into a
handful of geographic clusters~\cite{Wirtz2020}, while full-graph models resolve
the network node by node from GIS data~\cite{vanderHeijde2017}, an approach demonstrated in
detailed dynamic-simulation libraries~\cite{Giraud2015}. Single-node
representations remain standard in national-scale
analysis~\cite{Connolly2014_EnergyPLAN}, and \citet{Gabrielli2018} likewise reduce the
thermal network of a distributed multi-energy system to a single aggregate efficiency.

Where the spectrum has been traversed, the comparison has typically confounded topology
with losses: moving from a lumped to a resolved network introduces pipe losses at the same
time as spatial structure, so a reported cost difference cannot be attributed to either.
This confound is structural rather than incidental: the objection that a coarse model
``misses losses'' is empirically indistinguishable from ``misses topology''. Recent full-graph studies do not resolve it,
because they do not vary resolution at all: \citet{Trentmann2026} embed full-graph
network expansion in an energy-hub MILP with temporally resolved heat-pump performance,
\citet{Ceruti2025} formulate a full-graph design MILP with storage and conclude that
temporal rather than spatial resolution dominates design quality, and
\citet{Vieth2026} treat expansion planning of existing full-graph networks under
resilience constraints. Each resolves the network fully throughout, so the aggregated
comparison studied here lies outside their scope. Closest to the decision-oriented view
taken here, \citet{Friedrich2025} trace how the fidelity of the network model propagates
into operational efficiency; but they too assess the simulated operation of a single model
rather than the regret of executing a lower-fidelity schedule under a higher-fidelity one,
and do not separate loss visibility from spatial topology.

Beyond topology, pipe models differ substantially in thermo-hydraulic fidelity.
Steady-state thermal losses follow from the fluid--ground temperature difference and the
pipe insulation, and are MILP-compatible whenever temperatures enter as fixed
parameters~\cite{Frederiksen2013}. Pressure drop is not: the Darcy--Weisbach relation is
quadratic in flow velocity and pumping power adds a cubic term, handled in the literature
either by piecewise linearisation~\cite{Bordin2016} or by second-order cone
relaxation~\cite{Blommaert2020}; \citet{Hari2024} integrate Darcy--Weisbach pressure drop
with explicit pump placement in a general thermal network flow framework. Reported pumping
costs span 1--5\,\% of thermal energy cost~\cite{Frederiksen2013}. This paper returns to
that range, since the network studied here falls two orders of magnitude below it for
reasons that are physical rather than numerical. Temperature propagation along a pipe
introduces exponential and bilinear terms~\cite{vanderHeijde2017}, validated by
\citet{Maldonado2024} against industrial data to a supply-temperature mean absolute error
below 0.5\,K, and transport delays of one to three hours arise in long
pipes~\cite{Giraud2015}. Across dispatch formulations, methods span linear, mixed-integer
and nonlinear programming~\cite{Ommen2016}; MIQCP occupies a middle ground in which
quadratic constraints are handled natively while integer variables retain
unit-commitment logic. \citet{Hering2021} demonstrate the approach for design and
dispatch with multiple heat pumps at the scale of tens of consumers, but report no
reference comparison isolating the cost of the piecewise-linear approximation against the
native quadratic formulation. The magnitude of that linearisation error in a
district-heating dispatch setting has not been published.

Two features of this literature limit what its comparisons can conclude, and both are
visible in Table~\ref{tab:priorwork}. The first is the confound above: no prior study
separates the effect of representing thermal losses from the effect of representing
spatial structure, so the two are never independently priced. The second is more
fundamental. Fidelity is judged by comparing \emph{objective values} across formulations,
which measures the models against one another; but a scheduling model is used to make a
decision, and what matters is the cost of having decided with the simpler model and then
operated the network. A formulation can misstate its objective yet produce a schedule
that executes well, or state its objective plausibly yet schedule badly, and an
objective-difference comparison cannot distinguish the two. The closest prior work
approaches this without closing it: \citet{Quaggiotto2021} report for the Verona network
that simplifying thermal behaviour in a dispatch optimisation underestimates operational
cost relative to a model-predictive controller using a detailed network model, and
\citet{Jansen2024} confirm the direction with a mixed-integer nonlinear MPC on a Belgian
network, but neither separates topology from thermo-hydraulic fidelity nor expresses the
gap as a function of observable network properties. \citet{Wirtz2021_complexity} vary the
level of detail across 24 MILP formulations of a multi-energy system and identify which
simplifications are cost-relevant, the controlled-comparison logic this paper adopts,
but their complexity axis spans technology and part-load representation rather than
network topology and thermo-hydraulics, and the building-scale system they study does not
resolve pipe-by-pipe structure. \citet{Kotzur2021} review complexity management for
energy-system optimisation more broadly and observe that, while temporal aggregation is
well studied, spatial aggregation ``has only a small community'' and lacks systematic
guidance.

A separate literature quantifies the dispatch value of the generation technologies
studied here while leaving the network lumped. \citet{Mollenhauer2018} formulate hourly
annual unit commitment for German CHP plants retrofitted with heat pumps and thermal
storage (methodologically close to the dispatch model used here, but representing the
network as a single aggregated demand node), and \citet{Walden2023} optimise an
industrial power-to-heat system with an accurate nonlinear component model and no network
representation at all. \citet{Nielsen2016} show for the Danish market that deterministic
dispatch systematically overestimates the economic value of heat-pump and electric-boiler
flexibility relative to stochastic formulations. \citet{Riedl2026} propose a flexibility
assessment framework for district heating and cooling that incorporates network-side
constraints alongside generation-side flexibility, though the published case studies do
not yet vary spatial resolution as an experimental factor. None of this work treats
network representation as a factor, so the interaction between topology and technology
value remains open.

This paper addresses both gaps directly: a control condition that separates loss
visibility from spatial resolution, and a decision-regret metric that separates
estimation error from decision quality. The second is not a new idea in optimisation:
it is analogous to the value of the stochastic solution~\citep{BirgeLouveaux2011}, transposed
from the uncertainty axis to the resolution axis. The analogy is structural rather than exact:
the forward valuation is an overlay under a common high-fidelity model, so, unlike the classical
value of the stochastic solution, it need not be sign-definite (a loss-blind schedule can
forward-evaluate as cheaper on a mispriced term). The claimed contribution lies in the
transposition and the empirical finding rather than in the concept itself.
